\documentclass[main.tex]{subfiles}
\begin{document}

\chapter{Protokol dan Manajemen Kunci}

\begin{subcpmk}
  \item Sub-CPMK 1.4: Mengenal peran protokol SSL/TLS dan PKI sebagai aplikasi algoritma modern
  \item Sub-CPMK 2.2: Menjelaskan urgensi kriptografi di TI lewat studi kasus HTTPS/TLS
\end{subcpmk}

\noindent\textit{Rujukan:} Munir \cite{munir_slides_itb}; Boneh \cite{boneh_cs255}; Stallings \cite{stallings2017}.

% ============================================================
% MATERI POKOK
% ============================================================
\input{bab/bab-09/section-9-1}
\input{bab/bab-09/section-9-2}
\input{bab/bab-09/section-9-3}

% ============================================================
% AKTIVITAS PEMBELAJARAN
% ============================================================

\begin{aktivitas}
  \item \textbf{Analisis}: Buka sertifikat situs web HTTPS (browser: klik ikon gembok $\to$ sertifikat). Identifikasi issuer, subject, chain of trust.

  \item \textbf{Studi Kasus}: Analisis bagaimana HTTPS (TLS) melindungi transaksi e-commerce. Jelaskan peran sertifikat.

  \item \textbf{Diskusi}: Apa risiko jika kunci privat CA root bocor? Apa itu certificate revocation (CRL/OCSP)?
\end{aktivitas}

% ============================================================
% LATIHAN DAN REFLEKSI
% ============================================================

\begin{latihan}
  \item Jelaskan chain of trust. Mengapa klien memerlukan daftar CA root tepercaya?

  \item Sebutkan fase-fase dalam siklus hidup kunci (generasi, distribusi, penyimpanan, rotasi, pencabutan).

  \item Bagaimana TLS menggabungkan kriptografi simetris dan asimetris? Apa peran Diffie-Hellman?

  \item Apa peran sertifikat digital dalam TLS? Apa yang terjadi jika sertifikat tidak divalidasi?
\end{latihan}

% ============================================================
% ASESMEN
% ============================================================

\begin{asesmen}
\textbf{Instrumen Penilaian untuk Sub-CPMK 1.4, 2.2}

Jelaskan dalam essay: (a) Peran PKI dan chain of trust dalam TLS; (b) Mengapa validasi sertifikat penting terhadap MITM; (c) Mengapa HTTPS penting bagi privasi dan transaksi elektronik.
\end{asesmen}

% ============================================================
% CHECKLIST KOMPETENSI
% ============================================================

\begin{checklist}
  \item Saya memahami PKI dan komponennya (CA, RA, revocation)
  \item Saya dapat menjelaskan chain of trust dan validasi sertifikat
  \item Saya memahami peran TLS dalam keamanan komunikasi (HTTPS)
  \item Saya dapat menganalisis siklus hidup manajemen kunci
\end{checklist}

% ============================================================
% RANGKUMAN
% ============================================================

\begin{rangkuman}
Bab ini membahas protokol keamanan (TLS/SSL, SSH, PGP) dan manajemen kunci (PKI, siklus hidup kunci, RNG). TLS menggabungkan semua komponen kriptografi dalam sistem komunikasi aman. PKI menyediakan infrastruktur untuk distribusi dan verifikasi kunci publik.

\textbf{Poin Kunci:}
\begin{itemize}
  \item TLS: handshake (RSA/DH), sertifikat X.509, cipher suite, record layer
  \item PKI: CA hierarchy, chain of trust, certificate revocation
  \item Manajemen kunci: generasi, distribusi, rotasi, pencabutan
  \item RNG: TRNG vs PRNG, keacakan untuk keamanan kriptografi
\end{itemize}

\textbf{Kata Kunci}: \crypt{TLS}, \crypt{SSL}, \crypt{PKI}, \crypt{Sertifikat Digital}, \crypt{Chain of Trust}, \crypt{Manajemen Kunci}, \crypt{RNG}
\end{rangkuman}

% ============================================================
% Persiapan untuk Bab Berikutnya
% ============================================================

\begin{transition}
Setelah mengenal PKI/TLS dan urgensi di TI, lanjut ke \textbf{Bab X: Evaluasi, Refleksi, dan Integrasi Kompetensi} (UTS/UAS sesuai CPMK-1 dan CPMK-2). Materi RNG/steganografi tersedia sebagai lampiran opsional.
\end{transition}

\ifSubfilesClassLoaded{
  \renewcommand{\bibname}{Daftar Pustaka}
  \bibliographystyle{plain}
  \bibliography{../references}
}{}

\end{document}
